% ID: SEL-D0014-BC-E05
% ESTADO: cotejada. Segunda pasada visual completa y borrador compilado revisados.
% Fuente: Selectividad/matematicas_aplicadas/matematicas_aplicadas_P_B.pdf
% SHA-256: 51601c4f5a6c1dfcd860f621a659f8b9c09f53ff4cd31a1fa1e827dbd8f06159
% Página: 2; bloque C; ejercicio original 5.

Se han mezclado 90 llaves electrónicas de apertura de un determinado
garaje, con apariencia idéntica, de las cuales 60 funcionan correctamente
y 30 no funcionan. Se eligen al azar 2 de las 90 llaves.
\begin{enumerate}[label=\alph*)]
  \item \textbf{(0.7 puntos)} \textquestiondown{}Cuál es la probabilidad de
  que las dos llaves elegidas abran la puerta del garaje?
  \item \textbf{(0.8 puntos)} \textquestiondown{}Cuál es la probabilidad de
  poder abrir el garaje con alguna de ellas?
  \item \textbf{(1 punto)} \textquestiondown{}Cuál es la probabilidad de que
  una de las llaves elegidas funcione correctamente y la otra no?
\end{enumerate}
